% cap6_concluzii.tex
\chapter{Concluzii și dezvoltări viitoare}
\label{ch:concluzii}

%=============================================================================
\section{Concluzii generale}
%=============================================================================

Lucrarea de față a propus și a implementat un sistem de casă inteligentă
\emph{local-first}, funcțional integral în rețeaua locală, fără dependență de
servicii cloud externe~\cite{kleppmann2019localfirst}.  Obiectivul central
formulat în \chapref{ch:introducere} — de a demonstra că securitatea,
confidențialitatea și autonomia unui mediu rezidențial pot fi asigurate prin
mijloace accesibile, rulate pe hardware de consum — a fost atins prin
construirea unui sistem stratificat, de la senzorii de câmp până la interfața
mobilă a utilizatorului.

Arhitectura realizată cuprinde cinci straturi funcționale: (1)~stratul
senzori și actuatori, compus din 11 piese hardware; (2)~stratul noduri
\gls{iot} pe bază de \mbox{ESP32}; (3)~gateway-ul Raspberry Pi
(\gls{rpi}), care găzduiește broker-ul \gls{mqtt} Mosquitto, baza de date
PostgreSQL/TimescaleDB și motorul de \gls{ml}; (4)~\gls{api}-ul REST
NestJS; și (5)~aplicația mobilă React Native.  Integrarea verticală a
acestor componente — firmware \(\leftrightarrow\) broker \(\leftrightarrow\)
\gls{api} \(\leftrightarrow\) aplicație mobilă — asigură un flux de date
continuu, de la citirea brută a senzorilor până la vizualizarea în timp real
și controlul actuatorilor.

Nodul interior, bazat pe modulul ESP32 WROOM-32, integrează doi senzori de
temperatură și umiditate \gls{dht}11, un detector de gaz MQ-2, două
fotorezistențe \gls{ldr} pentru măsurarea luminozității ambientale, trei
\gls{led}-uri, un buzzer, un servo SG90 pentru acționarea draperiilor și un
ventilator de 5\,V.  Camera exterioară ESP32-CAM (AI-Thinker, senzor
OV2640) furnizează flux video \gls{mjpeg} și detectează mișcarea prin
senzorul \gls{pir} conectat la GPIO13.  Toate aceste piese fizice, enumerate
detaliat în \chapref{ch:implementare}, formează stratul primar de percepție
și acțiune al sistemului.

Comunicația intra-rețea se realizează prin protocolul \gls{mqtt}
(\mbox{OASIS}~\cite{oasis2019mqtt}), cu topicuri structurate pe spațiu de
nume \texttt{smarthome/\{nodeId\}/...} și mecanisme de Last~Will pentru
detectarea deconectărilor.  Descoperirea automată a gateway-ului se face prin
\gls{mdns} (\texttt{smarthome.local}), eliminând necesitatea configurării
manuale a adreselor \gls{ip} în aplicația mobilă.

Contribuția originală a lucrării, detaliată în secțiunile dedicate din
\chapref{ch:metodologie} și \chapref{ch:rezultate}, constă în motorul de
detectare a anomaliilor comportamentale, implementat în Python pe
\gls{rpi}.  Motorul extrage 12 trăsături relevante din datele senzorilor,
include variabile temporale și contextuale, normalizează vectorii de
trăsături cu \texttt{StandardScaler} și aplică algoritmul Isolation
Forest~\cite{liu2008isolation}.  Inferența rulează local, în timp real, fără
transmiterea datelor personale în exterior, în acord cu principiile
\emph{local-first}~\cite{kleppmann2019localfirst} și cu cerințele
\gls{gdpr}~\cite{gdpr2016}.

%=============================================================================
\section{Contribuții}
%=============================================================================

Principalele contribuții aduse prin această lucrare sunt:

\begin{enumerate}

\item \textbf{Arhitectură local-first completă pe cinci straturi.}  Sistemul
demonstrează că un ecosistem de casă inteligentă poate fi construit integral
cu componente accesibile (ESP32, Raspberry Pi) și software open-source, fără
a recurge la platforme cloud proprietare.  Această abordare păstrează datele
personale în rețeaua locală și menține funcționalitatea deplină chiar în
absența accesului la internet.

\item \textbf{Motor de detectare a anomaliilor cu Isolation Forest, rulat la
marginea rețelei (\gls{edge}).}  Algoritmul Isolation Forest~\cite{liu2008isolation},
nesupervizat prin natură, se potrivește scenariului \gls{iot} în care
etichetele de anomalie nu sunt disponibile în prealabil.  Implementarea
locală pe \gls{rpi} elimină latența și dependența față de un server remote,
permi\-țând reacția în timp real la evenimente neobișnuite~\cite{shi2016edge,
cook2020anomalyiot}.  Trăsăturile inginerești — incluzând temperatura medie a
celor doi senzori \gls{dht}11, delta de temperatură pe o oră, numărul de
evenimente de mișcare în ultimele 10 minute și indicatori temporali — sporesc
capacitatea de discriminare a modelului față de o abordare bazată exclusiv pe
valorile brute ale senzorilor.

\item \textbf{Pipeline firmware–backend–mobil integrat.}  Firmware-ul
Arduino/C++ al nodurilor publică date la fiecare ciclu de eșantionare;
serviciul \texttt{mqtt.service.ts} din backend consumă aceste date, le
persistă în TimescaleDB și le furnizează aplicației mobile prin Socket.IO,
obținând o latență end-to-end sub o secundă în condiții de rețea locală.

\item \textbf{Interfață mobilă orientată spre siguranță.}  Aplicația React
Native expune un scor de siguranță calculat în timp real, o listă de alerte
cu severitate, istoricul grafic al senzorilor și controlul direct al
actuatorilor, oferind utilizatorului un punct unic de monitorizare și
intervenție.

\end{enumerate}

%=============================================================================
\section{Limitări ale implementării curente}
%=============================================================================

O evaluare onestă a sistemului impune recunoașterea limitărilor prezente:

\paragraph{Model neantrenat pe date reale de durată.}  Modelul Isolation
Forest a fost validat cu date sintetice și cu un set restrâns de înregistrări
reale colectate pe o perioadă scurtă.  Pragul \texttt{ANOMALY\_THRESHOLD =
-0{,}1} și parametrul \texttt{contamination = 0{,}05} au valori
ilustrative, stabilite prin inspecție inițială, nu prin optimizare pe un set
de date reprezentativ acoperind anotimpuri, pattern-uri de ocupare variate și
tipuri diverse de anomalii.  Performanța reală în producție (precizie,
acoperire, scor~F1) rămâne neevaluată riguros și constituie un obiectiv de
cercetare viitoare~\cite{davis2006prcurve,fawcett2006roc}.

\paragraph{Lipsa criptării pe canalul \gls{mqtt}.}  Broker-ul Mosquitto
funcționează pe portul 1883 fără \gls{tls}, ceea ce înseamnă că traficul de
date al senzorilor poate fi interceptat de orice dispozitiv aflat în aceeași
rețea locală.  Deși riscul este limitat la perimetrul \gls{lan}, absența
criptării reprezintă o vulnerabilitate față de atacatorii cu acces fizic sau
wireless la rețea~\cite{alrawi2019sokhome}.

\paragraph{Acoperire redusă a testelor automate.}  Validarea sistemului s-a
bazat preponderent pe teste manuale și pe modulul de simulare a datelor
(\texttt{SIMULATE=true}).  Testele unitare și de integrare acoperă parțial
serviciile backend; nu există teste de regresie automate pentru firmware și
nici teste end-to-end pentru fluxul complet senzor–aplicație mobilă.

\paragraph{Topologie cu un singur nod interior.}  Prototipul curent suportă
un nod interior și o cameră exterioară.  Adăugarea mai multor noduri în
camere diferite necesită extinderea schemei de topicuri \gls{mqtt}, a
modelului de date și a interfeței mobile, fără a fi testată în această
iterație.

%=============================================================================
\section{Direcții de dezvoltare viitoare}
%=============================================================================

Pe baza limitărilor identificate și a potențialului neexplorat al
arhitecturii, sunt propuse următoarele direcții de continuare:

\paragraph{Antrenarea și evaluarea riguroasă a modelului \gls{ml}.}
Colectarea unui set de date reale pe cel puțin 30–90 de zile, acoperind
comportamente normale și evenimente anomale etichetate manual, ar permite
optimizarea hiperparametrilor Isolation Forest și calculul metricilor
cantitative (precizie, recall, AUC-ROC)~\cite{davis2006prcurve,fawcett2006roc,
cook2020anomalyiot}.  Compararea cu alte metode nesupravizate (LOF~\cite{breunig2000lof},
One-Class SVM~\cite{scholkopf2001ocsvm}) ar valida alegerea algoritmică.

\paragraph{Securizarea comunicației \gls{mqtt} cu \gls{tls}.}  Configurarea
Mosquitto cu certificate X.509 și activarea autentificării mutual-TLS pe
portul 8883 ar elimina vulnerabilitatea de interceptare și ar alinia sistemul
la recomandările de securitate pentru dispozitivele \gls{iot}~\cite{alrawi2019sokhome}.

\paragraph{Suport multi-nod.}  Extinderea arhitecturii pentru a suporta mai
multe noduri ESP32 în camere diferite — cu agregarea datelor la nivel de
gateway și vizualizare per-cameră în aplicația mobilă — ar transforma
prototipul actual într-un sistem de monitorizare la scara unei locuințe
complete~\cite{stojkoska2017smarthome}.

\paragraph{Notificări push.}  Integrarea unui serviciu de notificări
(Firebase Cloud Messaging sau un serviciu auto-găzduit) ar permite alertarea
utilizatorului pe dispozitivul mobil chiar și atunci când aplicația este în
fundal sau telefonul nu este în rețeaua locală, crescând semnificativ
caracterul practic al sistemului.

\paragraph{Reguli de automatizare editabile din aplicație.}  Tabelul
\texttt{automation\_rules} există deja în schema bazei de date, iar \gls{api}-ul
expune endpointul \texttt{/api/rules}.  Implementarea unui editor vizual de
reguli în aplicația mobilă — de tipul „dacă temperatura depășește 30°C și
ora este între 10:00 și 22:00, pornește ventilatorul" — ar adăuga un nivel
semnificativ de personalizare, fără modificări de infrastructură.

\paragraph{Optimizări energetice.}  Adoptarea modului deep-sleep pe nodurile
ESP32 și trimiterea datelor la intervale configurabile (în loc de ciclu
continuu) ar reduce consumul energetic al nodurilor alimentate din baterie,
extinzând aplicabilitatea sistemului la scenarii de monitorizare în zone fără
priză permanentă~\cite{shi2016edge,premsankar2018edgeiot}.

\paragraph{Îmbunătățiri ale modelului \gls{ml} prin calcul la margine.}
Pe termen lung, migrarea motorului de inferență pe un accelerator dedicat
(de exemplu, Coral USB, Raspberry Pi AI Kit) sau explorarea unor modele mai
ușoare bazate pe rețele neuronale (autoencoder LSTM pentru serii de timp)
ar putea crește acuratețea detecției și reduce latența de inferență față de
implementarea actuală pe CPU~\cite{chalapathy2019deepanomaly,shi2016edge}.

%=============================================================================
\section{Remarci finale}
%=============================================================================

Sistemul realizat demonstrează viabilitatea unui ecosistem de casă
inteligentă construit integral pe principii \emph{local-first}, cu hardware
accesibil și componente software open-source.  Detecția de anomalii
comportamentale prin Isolation Forest, rulată local pe Raspberry Pi,
reprezintă un pas concret spre sisteme \gls{iot} rezidențiale care protejează
confidențialitatea utilizatorilor fără a sacrifica funcționalitatea.
Limitările identificate — model nevalidat pe date de durată, absența
criptării \gls{mqtt}, acoperire redusă a testelor automate — trasează o
agendă clară de îmbunătățire pentru iterațiile viitoare ale proiectului.
